\section{Synthesis and discussion}

\subsection{What is verified and what is validated}

The laminar case provides the strongest verification statement because both the local velocity profile and the global pressure gradient have analytical references. Their errors, $3.18\%$ and $1.24\%$, are consistent with a correct developed-flow solution and a discretisation error concentrated near the wall. Grid sensitivity below $1\%$ on the final refinement indicates that the remaining discrepancy is not dominated by gross under-resolution.

The turbulent case is a validation problem. No analytical velocity profile exists, and the RANS closure introduces model-form error. The standard $k$--$\varepsilon$ configuration provides the most balanced result for the present data: it has the lowest profile RMSE and a Darcy factor within $0.082\%$ of the experimental reference. The close friction-factor agreement should not be viewed in isolation because integral wall friction can agree even when local profile features differ. Using both metrics reduces that ambiguity.

The low-Reynolds-number closure requires a more cautious reading. Its four-mesh profile differences are non-monotonic, and its turbulence quantities continue to change between the available iteration exports. The friction-factor discrepancy can therefore not be attributed uniquely to closure behaviour or near-wall resolution from the present evidence. The limited spacing-exponent comparison shows small changes in the selected mean profile, but it does not close this verification gap.

The experimental analysis follows the same logic. Calibration establishes traceability from voltage to force before any fluid-mechanics coefficient is evaluated. PSV and force measurements then identify the same physical mechanism through different observables. The velocity field resolves the alternating wake spatially; the lift record resolves its integral loading signature.

For the PSV acquisition, a spatial scan of local spectra recovers the probe Strouhal number at most grid points that satisfy the data-availability criterion. Its peak distribution remains quantised by the finite-record frequency resolution, so the scan is used as a consistency check rather than evidence of a spatially varying shedding frequency. The load-cell slope also remains stable under single-point omission, while this check does not address settling or reference-mass uncertainty.

\subsection{Interpretation of the two Strouhal estimates}

The measured Strouhal numbers are 0.204 from PSV and 0.181 from force data. The absolute difference is 0.023. Three factors prevent treating it as a measurement inconsistency:
\begin{itemize}
    \item the bulk velocity doubles between the two acquisitions;
    \item the water depth changes from 0.42 to 0.45~m;
    \item the estimators use a local transverse velocity and a span-integrated force, respectively.
\end{itemize}
Both values remain representative of periodic circular-cylinder shedding. A stricter cross-validation would require synchronous PSV and force acquisition at the same operating point, with a common clock and documented frequency resolution.

\subsection{Limitations}

The following limitations define the scope of the conclusions:
\begin{itemize}
    \item pairwise grid differences are reported instead of a formal asymptotic GCI;
    \item the turbulent near-wall grid cannot be reconstructed completely from the compact exports;
    \item the uncertainty of the reference pipe-flow profiles is not included in the turbulence-model ranking;
    \item the PSV field contains missing vectors and the phase sequence is not ensemble averaged;
    \item the uncertainty budget is first order and excludes blockage, free-surface deformation, end effects, and alignment;
    \item the cylinder datasets correspond to different operating points and are not synchronous.
\end{itemize}

These restrictions define the range of the conclusions: agreement with the selected benchmarks and consistent identification of the dominant wake dynamics.
